Tras un estudio en cierta profundidad de los métodos comunmente utilizados para el objetivo que nos concierne ¡Se llega a la conclusión individual de que existen gran cantidad de métodos diversos! \cite{rafii2018overview}

En nuestro caso, y aprovechando el recurso en formato tutorial que presentan los autores anteriormente citados en esta sección, y tras una prueba inicial de generación de máscaras, se elige el modelo utilizado previamente por los autores cuyo funcionamiento apunta a la inferencia de máscaras tiempo-frecuencia para la separación de las fuentes.

A tener en cuenta:
\begin{itemize}
    \item El modelo no separa directamente onda temporal, si no que predice máscaras tiempo-frecuencia.
    \item La salida estimada es una magnitud que se calcula como máscara de salida por magnitud de la mezcla.
    \item En el entrenamiento se comparan magnitudes estimadas con magnitudes objetivo.
    \item En el proceso de inferencia de la máscara se reconstruye el audio usando la fase de la mezcla.
\end{itemize}
